\section{Symbolic representations.} \label{app:symbols}

% --- Helper Functions ---
\newcommand{\skillFigure}[3]{%
  % #1: Human-readable Skill Name (e.g. Finish)
  % #2: File type (e.g. precondition -> init, effect -> eff)
  % #3: symbol type (e.g. precondition, effect)
  \StrSubstitute{#1}{_}{ }[\skillname]
  \begin{figure}[htbp]
    \centering
    \begin{subfigure}[b]{0.9\textwidth}
      \centering
      \includegraphics[width=\linewidth]{pictures/#1_#2_icon.png}
      \caption{Iconic representation}
    \end{subfigure}\hfill
    \begin{subfigure}[b]{0.9\textwidth}
      \centering
      \includegraphics[width=\linewidth]{pictures/#1_#2_invariants.png}
      \caption{Categorical representation}
    \end{subfigure}
    \caption{'\skillname' skill #3 symbolic representation. Two types of visualising obtained symbol: iconic and categorical.}
  \end{figure}
}

\foreach \skill in {Start, Find_the_Door, Go_to_the_Door, Open_the_Door, Find_the_Goal, Go_to_the_Goal, Finish} {
    \skillFigure{\skill}{init}{precondition}
    \skillFigure{\skill}{eff}{effect}
}